\section*{الفصل \ref{ch:g11:mechanical-energy} --- الطاقة الميكانيكية وانحفاظها}

\begin{solution}{exo:g11:mechanical-energy:1}
(أ) $v = \qty{6.94}{m/s}$:
$E_k = \tfrac12 \times 90 \times 6.94^2 \approx \qty{2.2e3}{J}$.
(ب) $v = \qty{36.1}{m/s}$:
$E_k = \tfrac12 \times 1200 \times 36.1^2 \approx \qty{7.8e5}{J}$.
(ج) $\times 4$؛ $\times 9$ --- أي التربيع.
\end{solution}

\begin{solution}{exo:g11:mechanical-energy:2}
$\Delta E_p = 75 \times 9.81 \times 850 \approx \qty{6.3e5}{J}$. ولا:
\omterm{def:g10:relative-motion:frame}{فالمرجع} يزيح القيمتين بالقدر نفسه. ولا: فلا يدخل إلا فرق
الارتفاع.
\end{solution}

\begin{solution}{exo:g11:mechanical-energy:3}
$v = \sqrt{2 \times 9.81 \times 20} = \qty{19.8}{m/s} \approx
\qty{71}{km/h}$. و$mgh = \tfrac12 mv^2$: فتتلاشى الكتلة.
\end{solution}

\begin{solution}{exo:g11:mechanical-energy:4}
$W = \Delta E_k = 0 - \tfrac12 \times 1200 \times 25^2 =
\qty{-3.75e5}{J}$؛ $F = \num{3.75e5}/50 = \qty{7.5e3}{N}$.
\end{solution}

\begin{solution}{exo:g11:mechanical-energy:5}
$v = \sqrt{2 \times 9.81 \times 0.080} \approx \qty{1.3}{m/s}$؛ وترتفع
عائدةً إلى \qty{8.0}{cm} (بالانحفاظ).
\end{solution}

\begin{solution}{exo:g11:mechanical-energy:6}
$\num{1.5e9}/\num{2.6e3} \approx \num{5.9e5}$: أي نحو ستمئة
ألف رصاصة. و$h = 88.9^2/(2 \times 9.81) \approx \qty{4.0e2}{m}$
--- أي أن سرعته هو كانت ستقدر على رفع القطار $400$ مترًا.
\end{solution}

\begin{solution}{exo:g11:mechanical-energy:7}
السرعة نفسها $v = \sqrt{2 \times 9.81 \times 3.0} \approx \qty{7.7}{m/s}$
(فلا يدخل إلا الهبوط). أما الزحلوقة الحادّة فتفوز في الزمن: إذ تبلغ السرعة
العالية أبكر ومسارها أقصر.
\end{solution}

\begin{solution}{exo:g11:mechanical-energy:8}
$h = v^2/2g = 144/19.62 \approx \qty{7.3}{m}$. وعند \qty{4.0}{m}:
$v = \sqrt{144 - 2 \times 9.81 \times 4.0} \approx \qty{8.1}{m/s}$ ---
وهي نفسها في الاتجاهين: إذ إنّ $E_m$ تتوقف على الارتفاع لا على جهة الانتقال.
\end{solution}

\begin{solution}{exo:g11:mechanical-energy:9}
$v = \sqrt{23^2 + 2 \times 9.81 \times 35} = \sqrt{1216} \approx
\qty{35}{m/s}$. وهي أصغر في الواقع: \omterm{def:g11:circuits-and-power:resistance}{فمقاومة} الهواء تبدّد جزءًا من
$E_m$.
\end{solution}

\begin{solution}{exo:g11:mechanical-energy:10}
$\Delta E_m = \tfrac12 \times 45 \times 8.0^2 - 45 \times 9.81 \times
6.0 = 1440 - 2649 \approx \qty{-1.2e3}{J}$؛
$f = 1209/15 \approx \qty{81}{N}$. وتذهب حرارةً، في الزلاجتين وفي
الثلج.
\end{solution}

\begin{solution}{exo:g11:mechanical-energy:11}
$v = \sqrt{2 \times 9.81 \times 0.15} \approx \qty{1.7}{m/s}$؛ وترتفع
إلى \qty{15}{cm} مرة أخرى. فالوتد يغيّر \omterm{def:g10:relative-motion:trajectory}{المسار} (دائرة أضيق)، لا
\omterm{def:g10:energy-conservation:energy}{الطاقة}: فالارتفاعات والسرعات لا تُمسّ.
\end{solution}

\begin{solution}{exo:g11:mechanical-energy:12}
$v = \sqrt{2 \times 9.81 \times 2.5} = \qty{7.0}{m/s}$. وترتفع ارتفاعات الصعود
تناسبًا مع $E_m$، ومنه $\times 0.9$ في كل عبور: $2.25$، $2.03$، $1.82$،
$1.64$، $1.48$، $1.33$، $1.20$، $1.08$، $0.97$ --- فأول ما يقلّ عن
\qty{1.0}{m} هو العبور $9$.
\end{solution}

\begin{solution}{exo:g11:mechanical-energy:13}
\omterm{def:g10:energy-conservation:energy}{الطاقة} الابتدائية $E_m = \tfrac12 m \times 15^2 + m \times 9.81 \times 45$ نفسها
مهما تكن الزاوية، ومن ثم تكون سرعة الأرض
$v = \sqrt{225 + 2 \times 9.81 \times 45} = \sqrt{1108} \approx
\qty{33}{m/s}$. أما الزاوية فتثبّت \emph{اتجاه} سرعة
الهبوط، وزمن الطيران، والمدى --- لا سرعة الهبوط.
\end{solution}

\begin{solution}{exo:g11:mechanical-energy:14}
$\tfrac12 mv^2 = mgh - f\ell = 80 \times 9.81 \times 120 - 65 \times
800 = \num{94176} - \num{52000} = \qty{4.22e4}{J}$، ومنه
$v = \sqrt{2 \times \num{42176}/80} \approx \qty{32}{m/s}$
(\qty{117}{km/h}). والجزء المبدَّد: $52000/94176 \approx 55\%$.
\end{solution}

\begin{solution}{exo:g11:mechanical-energy:15}
$m = \qty{2.0e9}{kg}$:
$E_p = \num{2.0e9} \times 9.81 \times 300 \approx \qty{5.9e12}{J}
= \num{5.9e12}/\num{3.6e6} \approx \qty{1.6e6}{kW.h}$. والمسترجَع:
$0.85 \times \num{1.64e6} \approx \qty{1.4e6}{kW.h}$ --- أي نحو
$\num{1.4e5}$ بيت مدة يوم.
\end{solution}

\begin{solution}{pb:g11:mechanical-energy:1}
\textbf{1.} $v_{\text{القمة}} = \sqrt{gR} = \sqrt{9.81 \times 8.0}
\approx \qty{8.9}{m/s}$.

\textbf{2.} $E_m = mg\,(2R) + \tfrac12 m gR = \num{3.14e5} +
\num{7.8e4} \approx \qty{3.92e5}{J}$.

\textbf{3.} يعطي $mgH_{\min} = E_m$ النتيجة $H_{\min} = \qty{20}{m}$؛
وجبريًّا $H_{\min} = 2R + \frac{gR}{2g} = \frac{5R}{2}$.

\textbf{4.} $v_{\text{القمة}} = \sqrt{2g(H - 2R)} =
\sqrt{2 \times 9.81 \times 8.0} \approx \qty{12.5}{m/s}$؛
$v_{\text{القمة}}^2/(gR) = 2.0$ --- أي ضعف الحدّ الأدنى.

\textbf{5.} كل \omterm{def:g10:energy-conservation:energy}{طاقة} تتناسب مع $m$، ومن ثم تتلاشى $m$: فالسرعات نفسها
والأمان نفسه، ممتلئًا أو فارغًا.

\textbf{6.} $v = \sqrt{2 \times 9.81 \times 24} \approx \qty{21.7}{m/s}
\approx \qty{78}{km/h}$.

\textbf{7.} $v = \sqrt{2 \times 9.81 \times (24 - 8.0)} \approx
\qty{17.7}{m/s}$.

\textbf{8.} $mgH = \qty{4.71e5}{J}$. و$(E_p, E_k)$: عند الإطلاق
$(\num{4.71e5}, 0)$؛ وعند الأسفل $(0, \num{4.71e5})$؛ وعند الجانب
$(\num{1.57e5}, \num{3.14e5})$؛ وعند القمة $(\num{3.14e5}, \num{1.57e5})$ ---
ويجمع كل زوج إلى \qty{4.71e5}{J}.

\textbf{9.} $v = \sqrt{2 \times 9.81 \times (24 - 12)} \approx
\qty{15.3}{m/s}$.

\textbf{10.} $v(z) = \sqrt{2g(H - z)}$: فهو أسرع ما يكون عند أخفض نقطة من
\omterm{def:g10:relative-motion:trajectory}{المسار}، وأبطأ ما يكون عند قمة الحلقة.

\textbf{11.} الارتفاع نفسه، ومن ثم $\Delta E_p = 0$ و
$\Delta E_m = \Delta E_k = \tfrac12 \times 2000 \times (19.2^2 -
20.4^2) \approx \qty{-4.75e4}{J}$: فيجعل الوضع العجزَ حركيًّا
خالصًا، مقروءًا من مقياسي سرعة.

\textbf{12.} $W(\vect f) = \qty{-4.75e4}{J}$؛
$f = \num{4.75e4}/60 \approx \qty{7.9e2}{N}$.

\textbf{13.} \omterm{def:g10:energy-conservation:forms}{طاقة حرارية}، في العجلات والسكك والهواء.

\textbf{14.} $E_m$ عند الحسّاس الأول: $\num{4.16e5} + \num{3.9e4} =
\qty{4.55e5}{J}$؛ والضياع $\num{4.75e4}/\num{4.55e5} \approx 10\%$ لكل
\qty{60}{m}. وعند $H_{\min}$ تُبلَغ قمة الحلقة \emph{بالضبط}
عند $\sqrt{gR}$؛ فأي ضياع يهبط بها دون سرعة الأمان.

\textbf{15.} بضياع $\approx 10\%$ لكل \qty{60}{m}، يذهب نحو $15\%$ من
$mgH$ ($\approx \qty{7.1e4}{J}$): أي $E_m \approx \qty{4.00e5}{J}$،
وهي ما تزال فوق المطلوب \qty{3.92e5}{J} --- فآمن، ولا يفيض إلا $2\%$.
فالهامش لم يكن ترفًا.

\textbf{16.} $E_k = \tfrac12 \times 2000 \times 18.0^2 =
\qty{3.24e5}{J}$.

\textbf{17.} $F = \num{3.24e5}/45 = \qty{7.2e3}{N}$.

\textbf{18.} $a = 7200/2000 = \qty{3.6}{m/s^2} \approx 0.37\,g$:
أي حازم لكنه مريح.

\textbf{19.} \omterm{def:g10:inertia:inventory}{احتكاك} \omterm{def:g10:relative-motion:trajectory}{المسار}: $mgH - \num{3.24e5} = \num{470880} -
\num{324000} = \qty{1.47e5}{J}$. والتدقيق: $\num{1.47e5} + \num{3.24e5} =
\num{4.71e5} = mgH$ --- فكل \omterm{def:g11:work-of-force:work}{جول} محسوب.

\textbf{20.} ارتفاع الإطلاق $H = \qty{24}{m}$ ثبّت أمان
الحلقة، وكل سرعة $\sqrt{2g(H - z)}$، وفاتورة \omterm{def:g10:inertia:inventory}{الاحتكاك}، \omterm{def:g10:inertia:force}{وقوة}
الكبح؛ أما ثمن إهمال \omterm{def:g10:inertia:inventory}{الاحتكاك} فهو الهامش --- أي $\qty{4}{m}$
من التلّ التي يستهلكها العالم المحتكّ في صمت.
\end{solution}
